% Artículo sobre la ontología (perfil OHAC) del proyecto LAAD.
% Compilar con:  latexmk -pdf main.tex   (BibTeX + splncs04)
\documentclass[runningheads]{llncs}

\usepackage[T1]{fontenc}
\usepackage[utf8]{inputenc}
\usepackage[spanish,es-noshorthands,es-tabla,es-nodecimaldot]{babel}
\usepackage{microtype}
\usepackage{graphicx}
\usepackage{booktabs}
\usepackage{tabularx}
\usepackage{multirow}
\usepackage{xcolor}
\usepackage{listings}
\usepackage{tikz}
\usetikzlibrary{positioning,fit,arrows.meta,backgrounds}
\usepackage[hidelinks]{hyperref}
\usepackage{csquotes}

% Marcas de pendiente: se deben resolver antes del envío.
\newcommand{\pendiente}[1]{\textcolor{red}{[PENDIENTE: #1]}}
\newcommand{\ttt}[1]{\texttt{\small #1}}
\setlength{\emergencystretch}{3em}
\newcommand{\cita}[1]{«\textit{#1}»}

\lstdefinelanguage{turtle}{
  morekeywords={a,@prefix},
  sensitive=true,
  morecomment=[l]{\#},
  morestring=[b]",
}
\lstset{
  basicstyle=\ttfamily\scriptsize,
  columns=fullflexible,
  keepspaces=true,
  breaklines=true,
  frame=single,
  framerule=0.3pt,
  xleftmargin=2pt,
  xrightmargin=2pt,
  literate={á}{{\'a}}1 {é}{{\'e}}1 {í}{{\'i}}1 {ó}{{\'o}}1 {ú}{{\'u}}1 {ñ}{{\~n}}1 {«}{{\guillemotleft}}1 {»}{{\guillemotright}}1 {→}{{$\rightarrow$}}1,
}

\begin{document}

\title{Lo que se dijo, lo que consta y lo que se reconstruye:\\
un perfil de aplicación en cuatro capas para expedientes judiciales coloniales}
\titlerunning{Un perfil de aplicación en cuatro capas para expedientes coloniales}

\author{Jhon Jesus Quispe Machaca \and Ian Logan Will Quispe Ventura}
\authorrunning{J. J. Quispe Machaca e I. L. W. Quispe Ventura}
\institute{Universidad Nacional de San Antonio Abad del Cusco, Cusco, Perú\\
\email{215422@unsaac.edu.pe}, \email{211359@unsaac.edu.pe}}

\maketitle

\begin{abstract}
Los expedientes judiciales coloniales no registran hechos: registran lo que alguien afirmó, ante quién, en qué
actuación y con qué fuerza procesal. Los grafos de conocimiento que convierten esas afirmaciones en tripletas
«verdaderas», o que dan por hecho lo que resolvió el tribunal, borran justamente lo que interesa al historiador.
Presentamos OHAC, un perfil de aplicación para el Fondo Intendencia del Archivo Regional del Cusco que separa cuatro
capas en grafos nombrados: el documento (RiC-O), las menciones sin fusionar (PiCo y ROAR), las afirmaciones
atribuidas (CIDOC CRM E13 con el patrón STAR) y las reconstrucciones del investigador con autor, fecha y grado de
creencia (CRMinf), más una capa de procedencia del dato (PROV-O). Lo andino y lo procesal se expresan como un tesauro
SKOS y las reglas del perfil como formas SHACL. Un análisis de carencias muestra que el perfil no necesita declarar
ninguna clase ni propiedad propia. El caso piloto es una causa criminal de 1806 contra un cacique recaudador de
Tinta, con certificaciones de agosto desmentidas en declaraciones juradas de septiembre. Veinte preguntas de
competencia, redactadas antes de modelar, funcionan como pruebas automáticas; el grafo resultante es conforme a
SHACL, consistente bajo OWL~2~RL y responde 17 de las 20. La evaluación reveló además seis defectos de
interoperabilidad en los estándares reutilizados, entre ellos una presunción de sinceridad que CRMinf~1.2.1 impone a
toda afirmación de CIDOC CRM.

\keywords{Ontologías \and Archivos coloniales \and CIDOC CRM \and RiC-O \and Afirmaciones en disputa \and
Humanidades digitales \and Perú}
\end{abstract}

\noindent\textbf{Recurso.}
Tipo: perfil de aplicación (tesauro SKOS, formas SHACL, consultas SPARQL y grafo de ejemplo).
Licencia: \pendiente{licencia de datos y código}.
Repositorio: \pendiente{URL pública}.
DOI: \pendiente{DOI de Zenodo}.
Espacio de nombres: \ttt{https://w3id.org/laad/} \pendiente{registro en w3id.org}.

% =====================================================================
\section{Introducción}\label{sec:intro}

Un expediente judicial colonial es una cadena de actuaciones: peticiones, certificaciones, declaraciones juradas,
informes, vistas fiscales y autos. Cada pieza registra lo que alguien afirmó, ante quién, quién la escribió y quién
la firmó. Las afirmaciones se contradicen entre piezas, los testigos se retractan y el tribunal resuelve sin que su
resolución vuelva verdadero lo resuelto. Para el historiador, el dato es precisamente esa trama de voces.

Los grafos de conocimiento construidos sobre fuentes de este tipo suelen perder esa trama de tres maneras: convierten
cada afirmación en una tripleta sobre el mundo; toman la sentencia como el valor de verdad que zanja la disputa; y
fusionan en una sola persona todas las grafías de un nombre. Un prototipo anterior de nuestro equipo, construido
como grafo de propiedades a partir de la extracción de un modelo de lenguaje, cometió los tres errores: la vista
fiscal se convirtió en un estado «desmentido» de los hechos, la frase de un fiscal quedó atribuida a ocho
declarantes por una colisión de identificadores, y las grafías del denunciante se fusionaron sin evidencia
(sección~\ref{sec:requisitos}).

Este artículo presenta OHAC, un perfil de aplicación para el Fondo Intendencia del Archivo Regional del Cusco (ARC)
cuyo principio es separar lo que es el papel, lo que está escrito, lo que alguien afirmó y lo que un investigador
propone. El caso piloto es la causa criminal del Legajo~14, Expediente~18 (1806): Manuel Huancachoque, indio
originario del pueblo de Tinta, acusa al cacique recaudador Andrés Chuquitapa. En agosto, varios vecinos certifican
que Chuquitapa fue \cita{complice y coronel} de Túpac Amaru en 1780; en septiembre, ante el juez comisionado, uno de
ellos declara \cita{Comisionado por este y no coronel}, otro niega haber firmado y un tercero dice que no sabe
escribir. El subdelegado informa que las certificaciones las redactó el propio denunciante, y la vista fiscal
concluye que las quejas son \cita{menos conforme a la verdad}.

Las contribuciones son:
\begin{enumerate}
\item Un modelo en cuatro capas, más una de procedencia, separadas en grafos nombrados, que representa afirmaciones
  judiciales atribuidas sin valores de verdad y deja las identidades y los hechos para reconstrucciones con autor
  (sección~\ref{sec:modelo}).
\item Un análisis de carencias que muestra que ese modelo se expresa sin términos propios, reutilizando RiC-O~1.1,
  CIDOC CRM~7.1.3, CRMinf~1.2.1, PiCo, ROAR, PROV-O, SKOS y Web Annotation; lo local (lo andino y lo procesal) vive
  en un tesauro SKOS (sección~\ref{sec:carencias}).
\item Seis defectos de interoperabilidad en esos estándares, encontrados al validar y razonar sobre datos reales, con
  su tratamiento (sección~\ref{sec:defectos}).
\item Una evaluación reproducible: veinte preguntas de competencia redactadas antes de modelar y ejecutadas como
  pruebas, formas SHACL con pruebas de mutación, razonamiento OWL~2~RL y verificación literal de todas las citas
  (sección~\ref{sec:evaluacion}).
\end{enumerate}

% =====================================================================
\section{Trabajos relacionados}\label{sec:relacionados}

\paragraph{Descripción archivística.}
\emph{Records in Contexts} (RiC-CM~1.0 y su ontología RiC-O~1.1) es el modelo de referencia del Consejo
Internacional de Archivos~\cite{riccm10,rico11,clavaud2021rico}. Describe agrupaciones documentales, documentos e
instanciaciones, pero no la estructura argumentativa de lo escrito. EHRI construyó sobre RiC-O el mayor grafo de
archivos del Holocausto~\cite{garciagonzalez2023ehri}. CourtDocs propone un modelo para actas judiciales
históricas que combina RiC-O, PROV-O y CIDOC CRM~\cite{vafaie2023courtdocs}; a diferencia de nuestro perfil, usa
PROV-O para representar el propio proceso judicial, mientras que aquí PROV-O se reserva a la procedencia del dato y
el proceso se describe con actividades de CIDOC CRM. Las ontologías jurídicas generales, como LKIF
Core~\cite{hoekstra2009lkif}, modelan normas y argumentos del derecho moderno y no se ajustan a las actuaciones de
un pleito del Antiguo Régimen.

\paragraph{Afirmaciones sobre afirmaciones.}
El modelo de \emph{factoides} de la prosopografía~\cite{pasin2015factoid} registra cada dato junto con su fuente, pero
funde la fuente con quien afirma y no distingue roles procesales. FactGrid, sobre Wikibase, admite afirmaciones
contradictorias con calificadores como «según»~\cite{factgrid}; su exportación RDF, sin embargo, materializa como
tripleta directa la afirmación de mejor rango, que es lo que nuestras capas~1--3 prohíben. STAR reifica cada
afirmación como una asignación de atributo de CIDOC CRM (E13) con autoridad, fuente y evidencia~\cite{andrews2023star,andrews2025star}; es el patrón que adoptamos y extendemos.
CRMinf modela las creencias y las inferencias de quien argumenta~\cite{doerr2011factual,crminf121}. En el plano
genérico de RDF existen la reificación clásica, las relaciones n-arias~\cite{noy2006nary}, los grafos nombrados, las
nanopublicaciones~\cite{groth2010nanopub,kuhn2016nanopub} y los \emph{triple terms} de RDF~1.2, que en noviembre de
2026 siguen en fase de recomendación candidata~\cite{rdf12}. Massari et al.~\cite{massari2025} comparan esas
técnicas para datos culturales y concluyen que no hay una solución única; Mariani~\cite{mariani2023} observa que E13
se vuelve repetitivo para hipótesis compuestas y propone empaquetarlas como grafos, que es lo que hacemos en la capa
del investigador. Los grafos conjeturales~\cite{daquino2022conjectural} y su formalización
epistémica~\cite{vitali2026dec} resuelven un problema paralelo con predicados conjeturales propios.

\paragraph{Personas, lugares y fuentes coloniales.}
PiCo~\cite{pico} y ROAR~\cite{roar} separan la observación de una persona en una fuente de su reconstrucción. En
fuentes andinas, el conjunto de datos de Aucará~\cite{melo2026aucara} publica menciones de registros parroquiales
como insumo para el enlace probabilístico de registros, sin RDF. Candela et al.~\cite{candela2023colonial} proponen un
enfoque ontológico para el archivo colonial. Para lugares, el nomenclátor de HGIS de las
Indias~\cite{stangl2018hgis} cubre la América española de 1701 a 1808. No encontramos ningún grafo de conocimiento
del ARC ni de expedientes judiciales andinos con testimonios atribuidos; es una búsqueda negativa, no una prueba de
inexistencia.

% =====================================================================
\section{Requisitos y método}\label{sec:requisitos}

\paragraph{Metodología.}
Seguimos LOT~\cite{poveda2022lot} en sus fases de especificación, implementación, publicación y mantenimiento, con
el escenario de reutilización de recursos ontológicos de NeOn~\cite{suarez2015neon}. Los requisitos se escribieron
siguiendo la plantilla ORSD~\cite{suarez2009orsd} y se expresan como preguntas de competencia (PC)~\cite{gruninger1995,ren2014cq}.

\paragraph{Materiales.}
La fuente es una transcripción paleográfica del expediente, dividida en 21 tramos por los separadores del
transcriptor. La transcripción contiene solo cuatro marcas de folio y no indica recto ni vuelto; no se dispone de
las imágenes. Por eso la unidad citable es el \emph{tramo}, no el folio, y el folio consta solo cuando el texto lo
marca. Se descartó como fuente una versión modernizada del texto que resolvía variantes en silencio (fechas,
apellidos). La ficha del catálogo del ARC se \emph{contrasta} con la transcripción en lugar de asumirse: presenta al
menos cinco discrepancias, como una extensión de 19 folios incompatible con la marca \cita{[Folio 34]}.

\paragraph{Lecciones del prototipo previo.}
El diagnóstico del prototipo anterior (sección~\ref{sec:intro}) fijó requisitos no funcionales: ningún valor de
verdad en lo transcrito; ninguna fusión de personas sin una reconstrucción explícita; citas que sean subcadenas
exactas de la fuente; identificadores deterministas; y uso real de los estándares (sus IRI deben existir en las
versiones publicadas), no su mención en comentarios.

\paragraph{Preguntas de competencia.}
Se redactaron 20 PC \emph{antes} de modelar, como haría un historiador del Cusco colonial
(tabla~\ref{tab:pc}). Cada una indica la capa que la responde y la forma de su respuesta: dato documental (D), lista
de afirmaciones atribuidas (A), disputa («X lo afirma, Y lo niega»), \emph{no consta} (NC), reconstrucción con autor
(R) o procedencia (P). Seis PC tienen por respuesta correcta una disputa o un \emph{no consta}; es decir, el modelo
debe ser capaz de \emph{no} dar una respuesta factual. Las preguntas y unas respuestas de referencia provisionales,
cada una con su tramo, quedaron registradas en el control de versiones antes del primer cambio del modelo, lo que
funciona como un prerregistro.

\begin{table}[t]
\centering
\caption{Selección de las 20 preguntas de competencia. Capa: 1 documento, 2 menciones, 3 afirmaciones,
4 reconstrucción, P procedencia.}\label{tab:pc}
\scriptsize
\begin{tabularx}{\linewidth}{@{}lXcl@{}}
\toprule
PC & Pregunta & Capa & Forma \\
\midrule
01 & ¿Qué piezas forman el expediente, en qué orden, de qué tipo es cada actuación y qué fecha lleva? & 1+3 & D \\
05 & ¿Con qué grafías aparece el denunciante, en qué tramos y en qué pieza? & 2 & D \\
08 & Para cada testigo de agosto, ¿qué afirmó en su certificación y qué en su declaración de septiembre? & 3 & A \\
09 & ¿Fue Andrés Chuquitapa coronel de Túpac Amaru en la rebelión? & 3 & DISPUTA \\
10 & ¿Quién redactó las certificaciones de agosto? & 3 & DISPUTA \\
11 & ¿Cuánto tiempo estuvo preso Huancachoque? & 3 & DISPUTA \\
12 & ¿Se sacó finalmente a Huancachoque de la matrícula de tributarios? & 3 & NC \\
13 & ¿Qué castigo se ejecutó tras el auto del 26 de noviembre de 1806? & 3 & NC \\
14 & ¿A qué hora fue el alboroto ante la cárcel y cuántos indios participaron? & 3 & DISPUTA \\
16 & ¿Qué resolvió la autoridad y en qué afirmaciones previas se apoyó? & 3 & A \\
18 & ¿Qué cargos ejercía Chuquitapa, sobre qué ayllus y con qué vigencia conocida? & 3+4 & R \\
19 & ¿Son la misma persona el «Mariano Casorla» de septiembre y el «Martin Cazorla» de agosto? & 2+4 & R \\
20 & ¿De qué versión de la transcripción procede este dato, qué proceso lo extrajo y quién lo validó? & P & D \\
\bottomrule
\end{tabularx}
\end{table}

% =====================================================================
\section{El perfil OHAC}\label{sec:modelo}

\subsection{Principios}

\begin{enumerate}
\item \textbf{Capas separadas.} Cada capa vive en grafos nombrados propios: \ttt{g:documento},
  \ttt{g:texto/<lote>}, \ttt{g:afirmaciones/<lote>}, \ttt{g:reconstruccion/R-x} y \ttt{g:argumentacion/R-x}, más
  \ttt{g:procedencia}.
\item \textbf{Sin valores de verdad en las capas 1--3.} Una resolución es una afirmación con fuerza institucional
  resolutiva; no hace verdadero lo resuelto. Solo la capa~4 tiene creencias, siempre con autor y fecha.
\item \textbf{Sin fusión en la capa 2.} Cada aparición escrita de una persona, un grupo o un lugar es un nodo. Las
  identidades son hipótesis de la capa~4 y pueden competir entre sí.
\item \textbf{Lo local como datos.} Las condiciones andinas (originario, forastero, tributario), los cargos
  (cacique, recaudador, segunda), los tipos de actuación y los roles procesales son conceptos SKOS, no clases.
\item \textbf{Solo términos existentes} en las versiones publicadas de los estándares, comprobados contra sus
  archivos.
\item \textbf{PROV-O solo para la procedencia del dato}, nunca para \enquote{quién lo dijo en 1806}.
\item \textbf{Determinismo.} Un mapeo puro lleva registros planos al grafo; todo identificador se deriva del
  expediente, el tramo y un ordinal.
\end{enumerate}

\begin{figure}[t]
\centering
{\renewcommand{\ttt}[1]{\texttt{#1}}%
\begin{tikzpicture}[
  font=\tiny,
  nodo/.style={draw, rounded corners=2pt, fill=white, inner sep=2pt, align=center, minimum width=3.3cm},
  capa/.style={draw=gray!60, fill=#1, rounded corners=3pt, inner sep=4pt},
  flecha/.style={-{Stealth[length=4pt]}, thin},
  etq/.style={font=\tiny, fill=white, inner sep=1pt}
]
% capa 4
\node[nodo] (rec) at (0,4.6) {\ttt{picom:PersonReconstruction}\\ Andrés Chuquitapa};
\node[nodo] (inf) at (7,4.6) {\ttt{crminf:I5} (autor, fecha, argumento)\\ $\rightarrow$ \ttt{I2} creencia «probable»};
% capa 3
\node[nodo] (a1) at (0,3.0) {\ttt{E13} A-UT01-001 (agosto)\\ \emph{ejerce-cargo} $\rightarrow$ coronel rebelde\\ \ttt{pol:afirma}, \ttt{fza:testimonial}};
\node[nodo] (a2) at (7,3.0) {\ttt{E13} A-UT11-004 (septiembre)\\ \emph{ejerce-cargo} $\rightarrow$ coronel rebelde\\ \ttt{pol:niega}, \ttt{fza:testimonial}};
% capa 2
\node[nodo] (m1) at (0,1.5) {\ttt{PersonObservation}\\ «Andres Chuquitapa»};
\node[nodo] (m2) at (3.5,1.5) {\ttt{PersonObservation}\\ «Isidro Cano» (agosto)};
\node[nodo] (m3) at (7,1.5) {\ttt{PersonObservation}\\ «Isidro Cano» (septiembre)};
% capa 1
\node[nodo] (p1) at (0,0) {Pieza P01: certificación\\ \ttt{rico:Record}, \ttt{E65}};
\node[nodo] (p2) at (7,0) {Pieza P18: declaración jurada\\ \ttt{rico:Record}, \ttt{E65}};
\begin{scope}[on background layer]
\node[capa=violet!8, fit=(rec)(inf), label={[font=\scriptsize\bfseries]left:4}] {};
\node[capa=orange!10, fit=(a1)(a2), label={[font=\scriptsize\bfseries]left:3}] {};
\node[capa=green!8, fit=(m1)(m2)(m3), label={[font=\scriptsize\bfseries]left:2}] {};
\node[capa=blue!7, fit=(p1)(p2), label={[font=\scriptsize\bfseries]left:1}] {};
\end{scope}
\draw[flecha] (a1) -- node[etq,right]{P140} (m1);
\draw[flecha] (a1.south east) -- node[etq,pos=0.45]{P14} (m2);
\draw[flecha] (a2) -- node[etq,right]{P14} (m3);
\draw[flecha] (a2.west) -- node[etq,above]{P17 (responde a)} (a1.east);
\draw[flecha] (a1.south west) to[bend right=30] node[etq,left]{P9i} (p1.west);
\draw[flecha] (a2.south east) to[bend left=30] node[etq,right]{P9i} (p2.east);
\draw[flecha] (rec) -- node[etq,right]{\ttt{prov:wasDerivedFrom}} (a1);
\draw[flecha] (inf) -- node[etq,above]{\ttt{J4\_that} (grafo de la proposición)} (rec);
\end{tikzpicture}}
\caption{Las cuatro capas en el hilo central del caso. La afirmación de agosto y su negación de septiembre conviven
como dos E13 atribuidas; la identidad de Chuquitapa es una reconstrucción con autor y creencia. Se omiten las anclas
textuales (\ttt{oa:Annotation}) y la capa de procedencia.}\label{fig:capas}
\end{figure}

\subsection{Capa 1: el documento}
La jerarquía sigue al catálogo: fondo, sección, serie, legajo y expediente son \ttt{rico:RecordSet}. RiC-O tipifica
fondos, series y expedientes, pero no tiene \emph{sección} ni \emph{legajo}, que se expresan como conceptos del
tesauro. Cada pieza del expediente es a la vez un \ttt{rico:Record} y un \ttt{crm:E31\_Document}, creado por una
actuación (\ttt{crm:E65\_Creation}) con su tipo y su fuerza institucional. La capa~1 guarda solo la fecha \emph{tal
como está escrita}; una forma SHACL prohíbe la fecha normalizada en esta capa, porque normalizar \cita{1809} a 1805,
como hacía la versión modernizada, es ya una interpretación. La transcripción es una \ttt{rico:Instantiation}
derivada del original, identificada por su suma SHA-256, y cada tramo es un componente de ella. La extensión física
que declara el catálogo (19 folios) se registra con \ttt{rico:CarrierExtent}, con su fuente y la discrepancia anotada.

\subsection{Capa 2: las menciones}
Cada mención es una \ttt{picom:PersonObservation} (personas) o una \ttt{roar:Observation} (grupos, lugares y cosas)
con su nombre literal (\ttt{crm:E41\_Appellation}). Las menciones no se tipifican como \ttt{crm:E21\_Person}: hacerlo
afirmaría la existencia de la persona, y el propio expediente dice de uno de los certificantes, \cita{don Juan
Quiñones}, que \cita{no hay tal hombre en el pueblo de Tinta}. Toda mención lleva un ancla de Web
Annotation~\cite{webannotation} con un \ttt{TextQuoteSelector} (cita exacta, prefijo y sufijo) y un
\ttt{TextPositionSelector} sobre el tramo. Si quien es mencionado actúa en la pieza (firmante, receptor, declarante,
testigo de asistencia), el mapeo crea una participación n-aria \ttt{crm:PC14\_carried\_out\_by} con el rol
(\ttt{P14.1\_in\_the\_role\_of}) en la actuación. De esas participaciones salen, para cualquier afirmación hecha en
la pieza, las facetas \emph{ante quién}, \emph{quién firma} y \emph{quién escribe}.

\subsection{Capa 3: las afirmaciones}
Cada afirmación es una \ttt{crm:E13\_Attribute\_Assignment} según el patrón STAR~\cite{andrews2025star}
(listado~\ref{lst:e13}): quién afirma (\ttt{P14}, una mención), en qué actuación (\ttt{P9i\_forms\_part\_of}), sobre
qué (\ttt{P140}), qué atributo (\ttt{P177}, de una lista cerrada) y con qué valor (\ttt{P141}). La polaridad (afirma o
niega), el fundamento invocado (visto, oído, por voz pública) y la fuerza institucional se expresan con
\ttt{P2\_has\_type} hacia conceptos del tesauro. Las relaciones de discurso usan \ttt{P17\_was\_motivated\_by}: la
afirmación a la que se responde y la persona \emph{según} la cual se dice algo. Nos apartamos de STAR en tres puntos:
la fuente es la actuación más el ancla textual, no \ttt{P17}; \ttt{P14} es quien afirma en 1806 y no la autoridad
del investigador; y las facetas procesales (ante quién, quién firma) se derivan de las participaciones de la
actuación. Cuando el anotador resuelve un sujeto elíptico o distribuye una afirmación colectiva, la decisión se
marca con un concepto \ttt{anot:*} y una nota obligatoria.

\begin{lstlisting}[language=turtle, float=t, caption={La negación de Isidro Cano en septiembre (capa 3, abreviada).
La afirmación de agosto que niega es \ttt{A-UT01-001}.}, label=lst:e13]
id:afirmacion/A-UT11-004 a crm:E13_Attribute_Assignment ;
    crm:P14_carried_out_by           id:mencion/M-UT11-001 ;      # «Isidro Cano», declarante
    crm:P9i_forms_part_of            id:actuacion/P18 ;           # declaración jurada, septiembre
    crm:P140_assigned_attribute_to   id:mencion/M-UT01-004 ;      # «Andres Chuquitapa» (sujeto elíptico)
    crm:P177_assigned_property_of_type atr:ejerce-cargo ;
    crm:P141_assigned                cargo:coronel-rebelde ;
    crm:P2_has_type  pol:niega, fza:testimonial, fun:no-indica, anot:responde-a, anot:sujeto-eliptico ;
    crm:P17_was_motivated_by         id:afirmacion/A-UT01-001 .
id:anotacion/A-UT11-004 a oa:Annotation ; oa:hasBody id:afirmacion/A-UT11-004 ;
    oa:hasTarget [ oa:hasSource id:transcripcion/.../UT11 ;
                   oa:hasSelector [ a oa:TextQuoteSelector ; oa:exact "Comisionado por este y no coronel" ] ] .
\end{lstlisting}

\paragraph{Fuerza institucional.}
El tipo de actuación determina la fuerza de lo afirmado en ella: testimonial (certificación, declaración jurada,
ratificación), de parte (petición), informativa (informe, oficio), dictaminadora (vista fiscal), resolutiva (auto,
decreto) y fedataria (notificación, desglose). La fuerza es un dato del tesauro (\ttt{act:x skos:related fza:y}) y
no se escribe en el código. Distinguirla permite responder «¿qué resolvió la autoridad?» sin responder «¿qué
ocurrió?»: el informe del subdelegado es informativo aunque lo emita una autoridad, y además lo emite una parte
interesada.

\subsection{Capa 4: las reconstrucciones}
Cada reconstrucción es un par de grafos. El grafo de la proposición contiene la identidad
(\ttt{picom:PersonReconstruction}, \ttt{roar:LocationReconstruction}), el evento o el ejercicio de un cargo
(\ttt{crm:E7\_Activity} con el cargo como rol y su vigencia en un \ttt{E52}), derivados de menciones y afirmaciones
con \ttt{prov:wasDerivedFrom}. El grafo de argumentación, aparte, contiene la inferencia
(\ttt{crminf:I5\_Inference\_Making}, también \ttt{prov:Activity}) con autor, fecha, argumento y premisas, y la creencia
(\ttt{I2\_Belief}) cuya proposición (\ttt{J4\_that}) es el grafo anterior. Separar ambos grafos evita que la creencia
forme parte de la proposición que evalúa, que es también la solución de Mariani~\cite{mariani2023}. Las hipótesis
rivales coexisten: el declarante de septiembre es \enquote{probablemente} Martín Cazorla y \enquote{posiblemente}
Mariano Casorla, y ninguna se elige en silencio. Los lugares reconstruidos llevan el identificador de HGIS de las
Indias como \ttt{crm:E42\_Identifier}, con el conjunto de datos como fuente, y los de Wikidata o TGN como
identificadores adicionales y \ttt{rdfs:seeAlso}, nunca \ttt{owl:sameAs}; no se copian geometrías, cuya licencia
(CC BY-NC-SA) se transmitiría al grafo.

\subsection{Procedencia, tesauro y formas}
Cada grafo es un \ttt{prov:Entity} generado por una actividad que usó la transcripción (con su suma), asociada a un
agente (persona o programa) y a un plan (la guía o el \emph{prompt}). Las reconstrucciones llevan autor y fecha.

El tesauro tiene 12 esquemas y 164 conceptos: tipos de actuación, roles procesales, condiciones, cargos, tipos de
entidad, atributos, polaridad, fundamento, fuerza, creencia, niveles documentales y marcas del anotador. Conserva las
rarezas del documento sin asimilarlas: \cita{India Ordinaria} no se iguala a \emph{originaria}, y el ayllu es un
\emph{grupo}, no un lugar. Las alineaciones externas usan solo \ttt{skos:closeMatch} o \ttt{relatedMatch}, porque las
etiquetas engañan: en Wikidata, la etiqueta quechua \emph{ayllu} está asignada a los ítems \emph{familia} y
\emph{comunidad}, no al ítem del ayllu.

Dieciséis formas SHACL~\cite{shacl} hacen cumplir el perfil. Entre otras: toda afirmación tiene exactamente un
afirmante con un rol de la rama \emph{afirmante}, una actuación, un sujeto, un atributo de la lista cerrada, un valor,
una polaridad, un fundamento, una fuerza y un ancla, y \emph{ninguna} creencia; ningún concepto del tesauro se usa
como clase; toda reconstrucción tiene inferencia con autor, argumento, fecha y una creencia de la lista; y todo lote
de las capas~1--3 tiene actividad, agente, plan y transcripción.

% =====================================================================
\section{Análisis de carencias}\label{sec:carencias}

Para cada necesidad derivada de las PC se buscó primero un término en los estándares; solo lo inexistente se habría
añadido al perfil (tabla~\ref{tab:carencias}). El resultado es que OHAC declara \emph{cero} clases y \emph{cero}
propiedades: es un perfil compuesto por importaciones con versión fijada, el tesauro, las formas y alineaciones con
el tesauro de roles de PiCo. Una prueba automática comprueba que cada IRI usada de un estándar existe en la versión
descargada de ese estándar, y otra, que un pleito civil sintético entra en el modelo sin clases ni propiedades nuevas.
Lo que podría presentarse como extensión (polaridad, fundamento, fuerza, roles) son conceptos usados con
\ttt{P2\_has\_type} y \ttt{P14.1}, que es como CIDOC CRM prevé la especialización local mediante \ttt{E55\_Type}. El
coste es que un consumidor genérico de CIDOC CRM que ignore \ttt{P2\_has\_type} leerá «asignó \emph{coronel}» en una
negación; las formas SHACL obligan a declarar la polaridad y toda consulta la devuelve.

\begin{table}[t]
\centering
\caption{Análisis de carencias (extracto). Ninguna necesidad exigió un término propio.}\label{tab:carencias}
\scriptsize
\begin{tabularx}{\linewidth}{@{}Xlc@{}}
\toprule
Necesidad & Solución con términos existentes & ¿Término propio? \\
\midrule
Sección y legajo & \ttt{rico:RecordSet} + concepto \ttt{doc:seccion}/\ttt{doc:legajo} & No \\
Extensión en folios & \ttt{rico:CarrierExtent} + \ttt{rico:quantity} & No \\
Folio individual, recto/vuelto & \ttt{rico:Instantiation} + concepto \ttt{doc:folio} & No \\
Afirmación con actor, actuación y tiempo & \ttt{crm:E13} + \ttt{P9i} + \ttt{P14} & No \\
Ante quién, quién firma, quién escribe & \ttt{crm:PC14} + \ttt{P14.1} + rol del tesauro & No \\
Polaridad, fundamento, fuerza & \ttt{crm:P2\_has\_type} + concepto & No \\
«Responde a», «según» & \ttt{crm:P17\_was\_motivated\_by} + marca \ttt{anot:} & No \\
Mención tal como está escrita & \ttt{picom:PersonObservation}, \ttt{roar:Observation} & No \\
Cita exacta & \ttt{oa:TextQuoteSelector} + \ttt{oa:TextPositionSelector} & No \\
Creencia del investigador & \ttt{crminf:I5}, \ttt{I2}, \ttt{J4}, \ttt{J5} & No \\
Ejercicio de cargo con vigencia & \ttt{crm:E7} + \ttt{PC14} (cargo como rol) + \ttt{E52} & No \\
Procedencia del dato & PROV-O & No \\
\bottomrule
\end{tabularx}
\end{table}

\subsection{Defectos de interoperabilidad encontrados}\label{sec:defectos}

Usar los estándares con datos reales, validarlos y razonar sobre ellos reveló seis defectos
(tabla~\ref{tab:defectos}). No los corregimos en origen: los evitamos o los medimos, y se informan en cada ejecución
del razonador.

\begin{table}[t]
\centering
\caption{Defectos de interoperabilidad hallados en los estándares reutilizados.}\label{tab:defectos}
\scriptsize
\begin{tabularx}{\linewidth}{@{}llX@{}}
\toprule
Id & Estándar & Defecto y tratamiento \\
\midrule
D1 & ROAR 0.1 & Declara un dominio global (\ttt{roar:Reconstruction}) sobre \ttt{prov:wasDerivedFrom}. Como
\ttt{prov:hadPrimarySource} es subpropiedad de esta, el patrón de PiCo aplicado a una observación de ROAR es
inconsistente (el razonador lo detectó). Las observaciones de ROAR usan \ttt{roar:documentedIn}; quedan 59
tipificaciones espurias de observaciones de PiCo, que se informan. \\
D2 & PiCo 1.1.0 & El prefijo de roles declarado en la ontología no coincide con las IRI publicadas del tesauro de
roles. Se usan las IRI publicadas. \\
D3 & CRMinf 1.2.1 & \ttt{J2\_concluded\_that} es subpropiedad de relaciones temporales (P175, P185); una inferencia
RDFS ingenua añadiría relaciones temporales falsas. \\
D4 & CRMinf 1.2.1 & \ttt{I6\_Belief\_Value} no se declara como clase y el rango de \ttt{J5} es un literal; los valores
de creencia no pueden ser conceptos. Se usa la notación del concepto como literal, restringida por SHACL. \\
D5 & RiC-O 1.1 & Expresa la extensión en folios, pero no el folio individual direccionable, ni recto y verso, ni
\emph{sección} y \emph{legajo}; sus tipos documentales no incluyen tipos judiciales. Se resuelve con conceptos SKOS. \\
D6 & CRMinf 1.2.1 & Declara \ttt{E13} subclase de \ttt{I1\_Argumentation}, definida como argumentación
\enquote{honesta}, y lee toda E13 como una proposición que su autor cree verdadera (\ttt{J33}). Con CRMinf cargado,
cada testimonio hereda una presunción de sinceridad. Se documenta y se mide (74 afirmaciones; 12 negaciones). \\
\bottomrule
\end{tabularx}
\end{table}

El defecto D6 merece una discusión aparte. En el caso piloto, el subdelegado informa que el denunciante redactó las
certificaciones y las hizo firmar \cita{sin saber lo mas de sus contenidos}; en septiembre, un certificante declara
que no sabe escribir y que otro puso por escrito lo que él relataba, y otro, que no firmó la suya.
Atribuirles una creencia sincera en lo que firmaron es exactamente el juicio histórico que el modelo quiere dejar a
la capa~4. Interpretamos por eso cada E13 de la capa~3 como una \emph{afirmación profesada en el registro}, y no
como la expresión de una creencia; el razonador cuenta en cada ejecución cuántas afirmaciones reciben la lectura de
sinceridad y cuántas de ellas son negaciones. Este conflicto entre la semántica de CRMinf y el uso de E13 para
testimonios es, hasta donde sabemos, inédito, y afecta a cualquier proyecto que combine STAR con CRMinf.

% =====================================================================
\section{Evaluación}\label{sec:evaluacion}

La evaluación se apoya en comprobaciones que no dependen de anotadores externos. El grafo del hilo central, que
cubre la disputa sobre el pasado rebelde del cacique y las certificaciones desmentidas, se construyó a partir de 33
piezas, 66 menciones, 74 afirmaciones y 27 reconstrucciones (incluida la de las fechas normalizadas). Tiene 5\,746 tripletas en 59 grafos nombrados y usa 40
clases y 67 propiedades, todas de los estándares. La tabla~\ref{tab:resultados} resume los resultados; todas las
cifras se reproducen con las pruebas del repositorio (206 pruebas).

\begin{table}[t]
\centering
\caption{Resultados de la evaluación del perfil sobre el hilo central del caso.}\label{tab:resultados}
\scriptsize
\begin{tabularx}{\linewidth}{@{}Xll@{}}
\toprule
Criterio & Umbral & Resultado \\
\midrule
Citas que son subcadena exacta del tramo (única normalización: espacios) & 100\,\% & 140/140 \\
Conformidad SHACL del grafo & 0 violaciones & 0 \\
Mutaciones del grafo que SHACL debe rechazar & todas & 35/35 \\
Inconsistencias bajo OWL~2~RL (datos + estándares + perfil) & 0 & 0 \\
Trampas críticas de OOPS!~\cite{poveda2014oops} & 0 & 0 (1 importante, artefacto) \\
IRI de estándares inexistentes en su versión publicada & 0 & 0 \\
Clases o propiedades nuevas para un pleito civil sintético & 0 & 0 \\
PC respondidas por consultas con nombre & 20 & 17 (3 requieren más anotación) \\
PC de disputa o \emph{no consta} respondidas como tales & 6/6 & 6/6 \\
\bottomrule
\end{tabularx}
\end{table}

\paragraph{Preguntas de competencia como pruebas.}
Cada PC es una consulta SPARQL con nombre y una prueba que compara su resultado con la respuesta de referencia. Las
tres PC pendientes (dos a escala de serie, que requieren extraer las fichas del catálogo, y la cadena completa de
autoridades) no exigen cambios del modelo, sino más anotación. Las disputas no se almacenan: se calculan agrupando
las afirmaciones por sujeto y atributo y comparando polaridades y valores. Para la PC09, el grafo devuelve que
Cano, en su certificación de agosto, y Orue, en su certificación y en su declaración, afirman que Chuquitapa fue
coronel; que Cano lo niega en su declaración jurada de septiembre; y que ninguna autoridad se pronuncia sobre el
punto. Cada afirmación lleva la cita literal, el tramo, el folio (o \emph{no consta}), la fuerza institucional y
quién recibió y firmó la pieza. Sobre el hilo central se calculan ocho disputas, más la de la firma de una de las
certificaciones, que su supuesto firmante niega.

\paragraph{Formas y mutaciones.}
Una forma que nunca falla no prueba nada. Por eso cada forma tiene al menos una prueba negativa: se altera el grafo
válido (se quita la polaridad de una afirmación, se añade una creencia en la capa~3, se usa un concepto como clase,
se normaliza una fecha en la capa~1, etc.) y se exige que SHACL lo rechace. Esta práctica reveló una forma que
aprobaba en vacío porque su objetivo no seleccionaba ningún nodo.

\paragraph{Razonamiento.}
El cierre OWL~2~RL de la unión de datos, perfil y estándares (unas 30\,600 tripletas, 104\,000 tras el cierre) no
produce inconsistencias. Fue el razonador el que detectó el defecto D1, al inferir una instancia de dos clases
disjuntas de ROAR.

\paragraph{Revisión independiente.}
Un revisor sin acceso a la documentación de diseño examinó el código, los datos y los estándares. Señaló tres
problemas críticos (los cargos con vigencia no estaban implementados; el código suponía un único expediente; faltaban
las facetas de quién escribe y el discurso referido) y ocho importantes, entre ellos que la capa de menciones
afirmaba la existencia de las personas. Todos se resolvieron o se documentaron como límites antes de esta versión.

% =====================================================================
\section{Discusión y limitaciones}\label{sec:discusion}

\paragraph{RDF frente a un grafo de propiedades.}
Elegimos RDF con grafos nombrados porque los términos de los estándares son los propios datos, la separación de
capas lleva procedencia por grafo y SHACL y OWL~RL permiten validar y razonar. El grafo de propiedades representa
con más naturalidad los atributos de una arista, pero en nuestro prototipo previo los tipos de relación libres
sustituyeron a los estándares. Un grafo de propiedades puede generarse como vista de lectura a partir del TriG, nunca
como fuente.

\paragraph{Negación.}
La polaridad se expresa como un tipo de la E13. La alternativa sin términos propios sería que la negación fuese una
E13 cuyo sujeto es la afirmación negada; solo es posible cuando la afirmación negada consta en el expediente (4 de
las 12 negaciones), por lo que se mantuvo la forma uniforme.

\paragraph{Limitaciones.}
(i)~El caso es un solo expediente, del que se anotó el hilo central; un segundo expediente, sorteado a ciegas con
semilla fija entre 3\,651, servirá para comprobar la generalización cuando se disponga de su transcripción.
(ii)~Las anotaciones del caso no han sido validadas por un especialista independiente: un solo autor actúa como
validador y fijó las respuestas de referencia, por lo que no se informa acuerdo entre anotadores.
(iii)~Sin imágenes, el folio y el lado de cada tramo no constan, salvo por las cuatro marcas del texto.
(iv)~Unos cien conceptos del tesauro carecen aún de definición.
(v)~OWL~2~RL es incompleto y CIDOC CRM se distribuye en RDFS, de modo que la ausencia de inconsistencias es una
condición necesaria, no suficiente.

% =====================================================================
\section{Conclusiones y trabajo futuro}\label{sec:conclusiones}

OHAC muestra que un expediente judicial colonial puede representarse sin decidir en los datos quién decía la verdad:
la trama de voces queda en afirmaciones atribuidas con fuerza institucional, y lo que el investigador propone queda
en reconstrucciones con autor, fecha y creencia. Todo ello cabe en los estándares existentes, cuyo uso con datos
reales reveló seis defectos que comunicaremos a sus mantenedores. El trabajo siguiente es llenar el modelo a escala:
un modelo de lenguaje rellenará registros planos con cita literal obligatoria y un mapeo determinista los llevará al
grafo, de modo que el modelo de lenguaje no decida identidades, verdades ni folios. Ese trabajo, y su evaluación
frente a líneas base sin grafo, quedan fuera de este artículo.

\paragraph{Disponibilidad del recurso.}
El perfil (tesauro, formas, consultas), el mapeo, el grafo de ejemplo, las pruebas y la documentación de diseño
(incluido el registro de decisiones) se publican en \pendiente{repositorio y DOI}, bajo \pendiente{licencias}.

\paragraph{Uso de inteligencia artificial generativa.}
\pendiente{Redactar con los autores.} El modelado, el código y las anotaciones del caso piloto se elaboraron con
asistencia de un modelo de lenguaje bajo la dirección de los autores; las anotaciones no han sido validadas por un
especialista independiente.

\bibliographystyle{splncs04}
\bibliography{referencias}

\end{document}
